\section*{第 \ref{ch:g8:protecting-the-nervous-system} 章——保护神经系统}

\begin{solution}{exo:g8:protecting-the-nervous-system:1}
它巩固当天增强的\omterm{def:g8:nervous-communication:synapse}{突触}——
把所学内容“归档”——并负责脑内的化学“清理”。
青少年的\omterm{def:g5:brain-in-command:system}{大脑}最需要\omterm{prop:g1:healthy-habits:needs}{睡眠}：
此时正处于\omterm{def:g8:nervous-communication:neuron}{神经}连接重组期，
\omterm{def:g6:exploring-our-environment:components}{生物}钟又因早起上学而推迟。
\end{solution}

\begin{solution}{exo:g8:protecting-the-nervous-system:2}
因为耳朵的感觉感受器会因剂量而死亡——
剂量即音量乘以暴露时间——且永不再生。
耗尽的是感受器这份“资本”。警示信号是：
喧闹的夜晚之后出现的耳鸣。
\end{solution}

\begin{solution}{exo:g8:protecting-the-nervous-system:3}
作用于\omterm{def:g8:nervous-communication:synapse}{突触}——模拟信使、阻断信使，
或迫使其释放。
\end{solution}

\begin{solution}{exo:g8:protecting-the-nervous-system:4}
它使各处\omterm{def:g8:nervous-communication:synapse}{突触}普遍受到阻尼：判断（权衡过程）、
\omterm{def:g5:brain-in-command:reaction}{反应时间}（回路的“时钟”）、平衡与记忆——
在每个环节都变得更慢、更迟钝。
\end{solution}

\begin{solution}{exo:g8:protecting-the-nervous-system:5}
反复摄入会使奖赏\omterm{def:g8:nervous-communication:synapse}{突触}产生\omterm{def:g4:adapted-to-their-world:adaptation}{适应}：
受体减少，阈值升高，
直到个体需要该物质才能感到正常。缺失它时，
\omterm{def:g5:brain-in-command:system}{神经系统}随之“抗议”——脱离\omterm{prop:g5:healthy-choices:substances}{成瘾}
需要逆着已重塑的\omterm{def:g8:nervous-communication:synapse}{突触}爬坡。
\end{solution}

\begin{solution}{exo:g8:protecting-the-nervous-system:6}
因为它训练的是同样的奖赏\omterm{def:g8:nervous-communication:synapse}{突触}，无需任何分子：
不可预测的小奖励，是最有力地强化\omterm{def:g8:nervous-communication:neuron}{神经}连接的
方式，
而其代价落在本章所要保护的\omterm{prop:g1:healthy-habits:needs}{睡眠}与注意力上。
\end{solution}

\begin{solution}{exo:g8:protecting-the-nervous-system:7}
复习强化\omterm{def:g8:nervous-communication:synapse}{突触}，\omterm{prop:g1:healthy-habits:needs}{睡眠}则巩固\omterm{def:g8:nervous-communication:synapse}{突触}。
凌晨零点至两点才学习的学生，只换来了生硬的连接，
却丢掉了归档；十点入睡的学生两者兼得。
未经巩固的复习，只有一半是自己的。
\end{solution}

\begin{solution}{exo:g8:protecting-the-nervous-system:8}
青少年的\omterm{def:g5:brain-in-command:system}{大脑}仍在“施工”中——紊乱的指令
会在整个建造过程中回响（见\cref{rem:g8:puberty:moods}中的\omterm{def:g8:nervous-communication:neuron}{神经}重组），
而且其\omterm{def:g8:nervous-communication:synapse}{突触}最具可塑性，
因此\omterm{prop:g5:healthy-choices:substances}{成瘾}所引起的奖赏回路重塑，
恰恰在此处最快发生。
\end{solution}

\begin{solution}{exo:g8:protecting-the-nervous-system:9}
汽车在踩下刹车前行驶的米数，正是\omterm{def:g5:brain-in-command:reaction}{反应时间}折算成的距离——
酒精会延长这一回路，使其\omterm{def:g8:nervous-communication:synapse}{突触}受到阻尼：
报告、权衡与执行，都变得更慢。
从机制上看，法律规定的限速，
实际上是对额外行驶距离的限制。
\end{solution}

\begin{solution}{exo:g8:protecting-the-nervous-system:10}
睡足时长：巩固在夜间进行。
控制响度：感受器是不可再生的资本。
避开\omterm{def:g8:nervous-communication:synapse}{突触}欺骗者：这些物质使可塑的\omterm{def:g8:nervous-communication:synapse}{突触}
朝陷阱方向重塑，而仍在发育的\omterm{def:g5:brain-in-command:system}{大脑}重塑得最快。
（关于头盔：\omterm{def:g8:nervous-communication:neuron}{神经}连接不像\omterm{def:g3:bones-and-muscles:skeleton}{骨骼}那样会自行愈合；
请善用这台“机器”：
\omterm{def:g8:nervous-communication:synapse}{突触}随使用而强化。）
\end{solution}

\begin{solution}{exo:g8:protecting-the-nervous-system:11}
寻求帮助成为一种日常维护，而非“坦白”：
\omterm{def:g8:nervous-communication:neuron}{神经}连接问题有治疗方法，可以求助，也无需羞耻——
因此陷入困境的人会更早寻求帮助。
提供帮助也变得切实可行：朋友的问题，
是可以陪同其走向护士办公室的事，
而不是需要保守的秘密，或需要下定论的事。
\end{solution}

\begin{solution}{exo:g8:protecting-the-nervous-system:12}
可调节的间隙正是机器能够学习的原因：使用
会强化连接，技能、记忆以及每一个经过训练的回路，
都是被巩固下来的\omterm{def:g8:nervous-communication:synapse}{突触}——
正是杂耍者一个月练习的积累。同样的可调节性
也正是脆弱之处：能对奖赏\omterm{def:g8:nervous-communication:synapse}{突触}施加剂量的物质与奖励安排，
同样会调节它们，使人需要
该剂量——\omterm{prop:g5:healthy-choices:substances}{成瘾}正是学习机制被教错了方向。
同一属性，两面作用：保护可调节性，
它会累积成技能；
若交给\omterm{def:g8:nervous-communication:synapse}{突触}欺骗者，它便累积成陷阱。
\end{solution}

\begin{solution}{pb:g8:protecting-the-nervous-system:1}
\textbf{1.} 预测：休息充足的志愿者
能更快、各次尝试间也更稳定地命中厘米级目标。机制：
\omterm{prop:g1:healthy-habits:needs}{睡眠}巩固了当天的\omterm{def:g8:nervous-communication:synapse}{突触}，并清除了化学损耗——
\omterm{prop:g1:healthy-habits:needs}{睡眠}不足的回路，
在每个\omterm{def:g8:nervous-communication:synapse}{突触}处都反应迟缓。

\textbf{2.} “电表显示剂量率；耳塞能降低它。
耳朵的感觉感受器会因累积剂量而死亡，
且永不再生——听力是一种资本，
在演唱会和全音量耳机中耗损最大，
应以“更小声、更短时、戴耳塞”加以保护。”

\textbf{3.} 模仿：尼古丁——在奖赏\omterm{def:g8:nervous-communication:synapse}{突触}处冒充天然的
信使。阻断：致瘫毒液——\omterm{def:g3:bones-and-muscles:muscle}{肌肉}\omterm{def:g8:nervous-communication:synapse}{突触}的信使被拒之门外。
压制：酒精与示意图相符，充当全系统的阻尼器
（箭头：从头到尾抑制\omterm{def:g8:nervous-communication:synapse}{突触}间的
“对话”）。

\textbf{4.} 目标：同样的奖赏\omterm{def:g8:nervous-communication:synapse}{突触}，
同样受不可预测奖励程序的训练——
通知演示，就是用光和声音对分子进行的
给药演示。

\textbf{5.} 十四岁时，\omterm{def:g8:nervous-communication:synapse}{突触}仍在建造之中，
可塑性最强：同样的剂量，
对成年\omterm{def:g5:brain-in-command:system}{大脑}中已定型的\omterm{def:g8:nervous-communication:synapse}{突触}，改线更深、见效更快
——“会没事的”是按简单难度玩的。

\textbf{6.} 第一支香烟的代价，就是之后所有的香烟：
每次用药都会启动奖赏\omterm{def:g8:nervous-communication:synapse}{突触}的\omterm{def:g4:adapted-to-their-world:adaptation}{适应}，
而入口恰好只有一个决定那么宽
——出口却是多年的\omterm{def:g8:nervous-communication:neuron}{神经}重连与反抗。提供的
是陷阱的门，而不是一支香烟。

\textbf{7.} 电视按计划发放奖励；信息流则
不可预测地、一条一条地发放——
而不可预测的小奖励，恰恰是最能训练
\omterm{def:g8:nervous-communication:synapse}{突触}的程序，这就是刷手机的冲动
比节目更难停下来的原因。

\textbf{8.} 只靠停下来，就是在对抗已重连的
\omterm{def:g8:nervous-communication:synapse}{突触}，但没有人必须这样做：护士和医生会
对这种\omterm{def:g8:nervous-communication:neuron}{神经}状况保密处理——包括未成年人在内——
且再\omterm{def:g4:adapted-to-their-world:adaptation}{适应}是可以成功的。第一步：今天找一位
可信赖的成年人，因为沉默是唯一的败招。

\textbf{9.} “一个月的抛接球练习，会强化\omterm{def:g8:nervous-communication:synapse}{突触}，
直到三个球能自己飞起来——这就是学习，
而这正是你的\omterm{def:g8:nervous-communication:synapse}{突触}存在的意义。\omterm{prop:g5:healthy-choices:substances}{成瘾}的机制
完全相同：用药会强化错误的回路，
直到它自行运转。\omterm{def:g5:brain-in-command:system}{大脑}无法只给你其中一种，
而不冒另一种的风险——可塑性就是全部的礼物。
所以问题从来不是你的\omterm{def:g8:nervous-communication:synapse}{突触}会不会被训练，
而只是被训练成什么。”

\textbf{10.} 因为提出邀约的那一刻，是最糟糕的
权衡时机——压力正压在负责决策的
\omterm{def:g8:nervous-communication:synapse}{突触}上。提前平静地拟好一句话，
就能让拒绝变成反射：
预先加权的路由，在需要之前就已就绪。

\textbf{11.} 一个站得住脚的排序：先\omterm{prop:g1:healthy-habits:needs}{睡眠}（它维持其他
一切所获得的成效），然后不愚弄\omterm{def:g8:nervous-communication:synapse}{突触}的东西
（那是不可逆的陷阱），再是响度
（不可替代的资本），然后是屏幕时长，
最后是头盔。首选的理由：其他
各项的益处都在夜间得到巩固——\omterm{prop:g1:healthy-habits:needs}{睡眠}是其他一切所依赖的
维持。

\textbf{12.} “我们在保护可塑的\omterm{def:g8:nervous-communication:synapse}{突触}，它们将学会
你今后所知的一切——而十四岁
正处在这十年的中段，是它们最快、最易训练、
也最易陷入陷阱的年纪。”
\end{solution}
